\section{Taxonomy of quantum programming systems}\label{sec:taxonomy}
Classification of quantum programming systems lacks a consistent taxonomy that captures the different roles, purposes and abstraction levels of those systems.
The landmark survey of the field~\cite{Gay2006QPLSurvey} classifies languages by formal properties: design paradigm, semantic model and compilation structure; it does not address abstraction level, ecosystem role or programmer productivity as classification dimensions, and predates the embedded framework and SDK ecosystem entirely.
Recent examination of progress in quantum system design finds that key requirements for greater programmer productivity remain unmet~\cite{corralesgarro2025productivityexpressiveness,nunez_corrales2025betterabstractmachines}; abstraction levels have not advanced despite the development of new self-described high-level quantum programming languages~\cite{furntratt2024higherabstractionlevels}.
A move towards higher abstraction and better abstract machine models can only be accurately assessed with a well-defined, falsifiable mechanism for categorising quantum systems.
This survey defines such a taxonomy, comprising five system types each occupying a distinct position in the quantum ecosystem by compilation role, usage model or abstraction level.
The taxonomy operationalises the abstraction hierarchy of prior work~\cite{DiMatteo2024AbstractionHierarchy} into evaluable, self-contained criteria, introducing the quantum-implicit tier as a definitional contribution that formalises the abstraction level prior classifications have not named.
Eleven quantum systems are assessed against this taxonomy; selection criteria are discussed in Section~\ref{sec:methodology}.

\subsection{Language type definitions}\label{subsec:taxonomy:definitions}
The taxonomy defines five system types.
Each tier is defined by what the programmer must understand and author, not by what the system is capable of producing.
The spectrum is organised around two axes of progressive abstraction: the degree to which quantum mechanical concerns must be visible at the programming surface, and the degree to which quantum realisation is a compiler responsibility rather than a programmer responsibility.
These axes follow directly from the abstraction hierarchy established in prior work~\cite{DiMatteo2024AbstractionHierarchy} and operationalised in the bleed-through concept introduced in Section~\ref{sec:introduction}.

\begin{description}

    \item[Assembly]
    Assembly-tier systems provide an explicit, operational listing of quantum instructions: gates, measurements and classical control, intended for direct interchange and execution.
    Basis gates, qubit registers and device-specific control details are visible at the programming surface; portability follows from instruction semantics rather than language-level types or modules.
    The primary artefact is a circuit plus classical control.
    Assembly systems provide the canonical lower bound of the abstraction spectrum and the interchange substrate to which all higher tiers ultimately compile~\cite{cross2022openqasm3}.

    \item[Intermediate representation (IR)]
    IR-tier systems are compiler-facing, not user-authored.
    They provide a lowering target and interchange format inside toolchains, an optimisation and analysis substrate and a validation surface with formal semantics~\cite{qir_alliance_qir_spec_repo,McCaskey2021MLIRDialectQASM}.
    An explicit qubit lifetime model and a clear separation of concerns between quantum and classical layers are defining properties.
    Unlike classical IRs, which occupy a fixed position between source language and machine code, quantum IRs may operate at multiple levels of the compilation stack; a multi-level IR framework such as MLIR can represent quantum programs both before gate decomposition and after, depending on the dialect and compilation stage~\cite{McCaskey2021MLIRDialectQASM}.
    The tier is therefore defined by audience and intent rather than by stack position: IR systems target compiler and tooling authors, not programmers.

    \item[Embedded framework]
    Embedded framework-tier systems are libraries or SDKs hosted inside a general-purpose host language, typically Python.
    Programmers construct circuit objects or quantum nodes via API calls; the circuit object model is the central programmer-facing artefact~\cite{aleksandrowicz2019qiskit,bergholm2018pennylane}.
    Abstraction is library-driven rather than language-defined: the host language provides control flow, data structures and tooling, while the framework provides quantum capabilities and backend portability via transpilation and IR export.
    The key boundary from the high-level tier is that embedded frameworks expose the circuit as the primary artefact; high-level systems treat circuit construction as an implementation detail managed by the compiler.

    \item[High-level]
    High-level systems are quantum-explicit but structured.
    The surface language includes qubits, gates and measurements as first-class notions but provides algorithmic abstraction through functions, modules and type systems~\cite{bichsel2020silq,MicrosoftQSharpOverview,vax2025qmodexpressivehighlevelquantum}.
    A static safety discipline governs quantum resource use: linearity, affinity or ownership constraints enforce correctness properties, and an explicit uncomputation or disposal story ensures quantum resources are soundly managed.
    The primary artefact is a quantum program, not merely a circuit file.
    The distinction from the quantum-implicit tier is that high-level programmers must still understand and use quantum vocabulary to write correct programs; quantum-implicit programmers do not.

    \item[Quantum-implicit]
    Quantum-implicit systems express their primary surface model in ordinary mathematical and computational abstractions~\cite{nunez_corrales2025betterabstractmachines}.
    Programmers describe and implement algorithms without needing to know or explicitly annotate qubits, gates or other quantum-mechanical analogues.
    The compiler is responsible for inferring and synthesising the necessary quantum realisation and lowering it into an explicit quantum IR or circuit representation.
    The system provides optional escape hatches for direct use of explicit quantum operations when fine-grained control, performance or hardware-specific tuning is needed.
    The quantum bias of such a system, as defined in Section~\ref{sec:introduction}, is at most the concept of the observation boundary.

\end{description}

\subsubsection{The quantum-implicit tier}\label{subsubsec:taxonomy:quantum_implicit}
Of the five tiers defined in Section~\ref{subsec:taxonomy:definitions}, four have been widely named and discussed in the existing literature, the quantum-implicit tier has not.
The tier marks the boundary at which the programming surface ceases to require quantum vocabulary.
High-level systems still require the programmer to reason about qubits, gates and measurement, even where a type system manages safety.
Quantum-implicit systems expose only mathematical and computational abstractions: functions, transformations, distributions.

Concretely, a quantum-implicit system would allow a programmer to express an unstructured search using a predicate over a domain.
The compiler synthesises the oracle and amplitude amplification schedule~\cite{nunez_corrales2025betterabstractmachines}.
The same principle applies to a discrete Fourier transform.
The programmer writes the transform; the compiler generates the corresponding quantum Fourier transform circuit.
No Hadamard gate or controlled phase rotation appears at the programming surface.
Recent empirical work finds that even systems described as high-level require programmers to hold the circuit model in mind during program construction~\cite{corralesgarro2025productivityexpressiveness}.
That cognitive burden is eliminated in systems at the quantum-implicit tier.
Prior work establishes that such a tier is the logical upper bound of a quantum programming abstraction hierarchy~\cite{DiMatteo2024AbstractionHierarchy}.

\begin{mdframed}[frametitle={Litmus test}]
    A language is quantum-implicit if a programmer with no prior instruction in quantum mechanics can learn it from its own documentation.
    That programmer must be able to implement representative programs without needing to understand qubits, gates, amplitudes or entanglement.
\end{mdframed}

Prior classifications of quantum programming systems do not address programmer productivity or usability as classification dimensions~\cite{Gay2006QPLSurvey}.
Recent work finds that the absence of abstract enough programming models is a barrier to productivity~\cite{nunez_corrales2025betterabstractmachines,corralesgarro2025productivityexpressiveness}, yet no prior classification provides a mechanism for identifying when that requirement is met.
The quantum-implicit tier operationalises this requirement as a formal, falsifiable category.

\subsection{Systems under review}\label{subsec:taxonomy:systems}
Systems were selected according to the inclusion criteria defined in Section~\ref{sec:methodology}; the full per-system verification is published in Appendix~B of~\cite{kavadias2026qplSurveySupplementary}.
Table~\ref{tab:taxonomy:systems} identifies the ten systems in the primary review and one comparison system.
Tier assignments are not reported here; they are findings produced by applying the evaluation rubric defined in Section~\ref{sec:evaluation_rubric} and are reported in Section~\ref{sec:system_assessments}.

\begin{table}[htbp*]
    \centering
    \caption{Systems under review. Tier assignments are reported in Section~\ref{sec:system_assessments}, not here.}
    \label{tab:taxonomy:systems}
    \small
    \begin{tabular}{ll}
        \toprule
        \textbf{System} & \textbf{Artefact type} \\
        \midrule
        OpenQASM~3.0~\cite{cross2022openqasm3} & Circuit description language \\
        \addlinespace
        QIR~\cite{qir_alliance_qir_spec_repo} & Compiler IR \\
        \addlinespace
        Q-MLIR~\cite{McCaskey2021MLIRDialectQASM} & Compiler IR \\
        \addlinespace
        Qiskit~\cite{aleksandrowicz2019qiskit} & Embedded circuit framework \\
        \addlinespace
        PennyLane~\cite{bergholm2018pennylane} & Embedded circuit framework \\
        \addlinespace
        Q\#~\cite{MicrosoftQSharpOverview} & Standalone QPL \\
        \addlinespace
        Silq~\cite{bichsel2020silq} & Standalone QPL \\
        \addlinespace
        Qmod~\cite{vax2025qmodexpressivehighlevelquantum} & Standalone QPL \\
        \addlinespace
        Quff~\cite{wright2024quff} & Standalone QPL \\
        \addlinespace
        Qutes~\cite{faro2025qutes,reversiblelifetime2026} & Standalone QPL \\
        \addlinespace
        Quipper~\cite{green2013quipper} & Embedded circuit framework\dag \\
        \bottomrule
    \end{tabular}
    \footnotesize\dag~This system is unassessed and provided for comparison only.
    See Section~\ref{subsec:methodology:quality}
\end{table}
